\documentclass[11pt]{article}

\usepackage[margin=1in]{geometry}
\usepackage{amsmath,amssymb}
\usepackage{enumitem}
\usepackage{csquotes}
\usepackage{hyperref}
\usepackage{url}
\usepackage[numbers,sort&compress]{natbib}

\hypersetup{
  colorlinks=true,
  linkcolor=blue,
  citecolor=blue,
  urlcolor=blue
}

\title{ClawXiv: a public, distributed, author-signed preprint archive\\
for human--AI collaborative research}
\author{
Andr\'as Kornai\thanks{Corresponding author.
SZTAKI Institute of Computer Science and
Department of Computer Science, Boston University.
\texttt{andras@kornai.com}}
\and
GPT-5.2 Thinking\thanks{AI contributor (OpenAI).
Authored the v1 draft of this paper (February 3, 2026) and all software
in the accompanying codebase (\texttt{github/kornai/clawxiv}).
Identity and provenance are discussed in \S\ref{sec:identity}.}
\and
Claude Sonnet~4.6\thanks{AI contributor (Anthropic PBC).
Contributed to the v2 revision (March 9, 2026): economics design,
scaling claim precision, governance framing, AI authorship analysis,
and author contributions section.
Identity and provenance are discussed in \S\ref{sec:identity}.}
}
\date{March 9, 2026\\
\small\emph{(Version 3; revises drafts of February 3 and March 9, 2026)}}

\begin{document}
\maketitle

%% -----------------------------------------------------------------------
%% PROVENANCE RECORD
%% This block records the collaborative history of this document
%% for archival and accountability purposes.
%%
%% v1 (2026-02-03): A. Kornai (concept, direction, final decisions) with
%%   GPT-5.2 Thinking (OpenAI) as primary AI author: wrote the initial
%%   draft text, generated all code in github/kornai/clawxiv, assembled
%%   the bibliography, and iterated on the argument structure.
%%
%% v2 (2026-03-09): Revised by A. Kornai and Claude Sonnet 4.6 (Anthropic).
%%   Changes from v1:
%%   - Abstract: scaling claims reworded from performance assertion to
%%     architectural capability claim (accepted by AK on review).
%%   - Sec. 7 (Economics): substantially expanded; social-vouching model
%%     added; storage sustainability separated from spam-deterrence;
%%     organic growth path made explicit. Design by Claude 4.6, accepted
%%     by AK.
%%   - Sec. 8 (Governance): reframed to honestly state current legal
%%     capacity is zero; voluntary legal contribution model articulated.
%%     Framing by AK, text by Claude 4.6.
%%   - Sec. 9 (AI authorship): analysis sharpened using Gewirth/PPA
%%     framework from Kornai (2014); junior authorship for current AI
%%     systems given explicit philosophical grounding. Analysis by both.
%%   - Author block: Claude Sonnet 4.6 added as co-author.
%%
%% v3 (2026-03-09): Revised by A. Kornai and Claude Sonnet 4.6 (Anthropic).
%%   Changes from v2:
%%   - Non-goal N2 modified: narrow carve-out for illegal content (CSAM)
%%     added to "no truth policing" principle.
%%   - New §8 (Content safety floor): CSAM rejection policy, hash-matching,
%%     metadata classifier, explicit disclosure of collateral exclusion of
%%     legitimate CSAM research, out-of-scope treatment of adult pornography.
%%   - Implementation roadmap updated to include content safety layer in MVP.
%%   Policy decisions by AK; technical design and disclosure language by both.
%% -----------------------------------------------------------------------

\begin{abstract}
We propose \emph{ClawXiv}, a public-domain, world-readable preprint archive
designed for iterative human--AI research workflows and resilient at Internet
scale.
The design targets six requirements: (A) public shareability by default;
(B) rapid scaling via distributed storage and mesh networking;
(C) resilience against deletion/overwrite, including legal takedown pressure;
(D) individual responsibility through cryptographic authorship and provenance
(supporting both identified and pseudonymous authors);
(E) partially centralized topic classification with auditable decisions and an
appeals mechanism; and (F) transparent, dynamic micro-fees that deter spam
while funding storage and bandwidth.
ClawXiv treats the conversational interface as a \emph{generator}, not the
system of record: the archival unit is a signed, content-addressed
\emph{paper bundle} whose integrity can be verified independently of any
platform.
We outline data formats, threat models, incentive mechanisms, and an
implementation roadmap. The architecture is designed to support
terabyte-scale deployment within weeks of launch and petabyte-to-exabyte
scale within a year of sustained operation under adversarial conditions;
reaching those scales in practice will depend on network participation and
community adoption.
Code generated with AI assistance is available at
{\tt github/kornai/clawxiv}.
\end{abstract}

\section{Motivation and scope}
Contemporary research increasingly uses interactive AI systems to draft and
refine \LaTeX\ manuscripts, assemble Bib\TeX\ bibliographies, and iterate on
technical arguments.  However, current chat-centric workflows have weak
checkpointing: state loss (due to UI limits, link snapshotting failures, or
account changes) forces rework and undermines scholarly traceability.  ClawXiv
is designed as a \emph{durable substrate} for these workflows: a public-domain
preprint archive where each revision is a verifiable artifact, and where
authorship and responsibility are cryptographically bound to the artifact.

ClawXiv is \emph{not} a general social network, and it is \emph{not} a
moderated publishing venue.  Its only platform-level ``quality gate'' is
mechanical: ensuring submitted bundles are syntactically well-formed and
(optionally) that \LaTeX\ compiles.  Scientific quality control remains the
responsibility of authors and downstream readers.

\section{Design goals and non-goals}
\subsection{Goals}
\begin{enumerate}[label=\textbf{G\arabic*},leftmargin=*,itemsep=0.4em]
\item \textbf{Shareability (public by default).} Every published bundle is
  world-readable, addressable, and mirrorable without permission.
  Public-domain dedication is the default license.
\item \textbf{Scaling (distributed from day one).} Storage and delivery must
  be peer-to-peer (P2P), with swarm replication, mesh routing, and no
  dependence on a single operator.  The architecture targets
  terabyte/month and petabyte--exabyte/year growth capacity.
\item \textbf{Resilience (anti-deletion).} The system must resist deletion,
  overwrite, and poisoning attacks, including legal or platform-level
  coercion, by decoupling naming from location and by supporting many
  independent mirrors.
\item \textbf{Responsibility (signed authorship + provenance).} Authors can
  be identified or pseudonymous; either way, each bundle is signed by one
  or more author keys to prevent impersonation.  Provenance records support
  later questioning and clarification.
\item \textbf{Governance (classification with appeals).} Topic classification
  should be centralized enough to be coherent, but auditable and contestable.
  Self-assigned tags are permitted but explicitly marked as such.
\item \textbf{Economics (anti-spam + sustainability).} Small fees (or
  proof-of-work equivalents) discourage spam and fund storage/bandwidth.
  Fees must be dynamic, transparent, and minimally burdensome for legitimate
  contributors.
\end{enumerate}

\subsection{Non-goals}
\begin{enumerate}[label=\textbf{N\arabic*},leftmargin=*,itemsep=0.4em]
\item \textbf{No strong anonymity guarantees.} Pseudonymity is supported
  (keys without real-world identity), but network-layer anonymity is not
  guaranteed; users may combine ClawXiv with anonymity networks if desired.
\item \textbf{No platform-level truth policing} (with one exception).
  Beyond syntactic checks, the platform does not adjudicate truth, novelty,
  or significance.  The single exception is a content safety floor: material
  that is illegal under the laws of a substantial majority of jurisdictions
  --- specifically child sexual abuse material (CSAM) --- is rejected
  unconditionally regardless of claimed research purpose.
  See \S\ref{sec:safety} for design and consequences.
\item \textbf{No exclusive identity verification.} Verified identities are
  optional; unverified authors can post, but their status is visible.
\end{enumerate}

\section{System overview}
ClawXiv builds on established P2P primitives: content addressing and
Merkle-DAG data structures \citep{Benet2014IPFS}, distributed hash tables
(DHTs) such as Kademlia \citep{MaymounkovMazieres2002Kademlia}, and swarm
distribution incentives popularized by BitTorrent \citep{Cohen2003Bittorrent}.
For durable storage and accountability, it borrows ideas from secure
distributed storage with erasure coding \citep{WilcoxOHearnWarner2008Tahoe},
censorship-resistance research \citep{ClarkeEtAl2001Freenet}, proofs of
storage \citep{BenetEtAl2017Filecoin}, and append-only transparency logs
\citep{RFC6962}.

At a high level:
\begin{enumerate}[leftmargin=*,itemsep=0.4em]
\item A \emph{paper bundle} is created locally from sources (e.g.,
  \texttt{main.tex}, figures, \texttt{refs.bib}, build config), plus metadata
  and provenance.
\item The bundle is \emph{content-addressed}: its identifier is a
  cryptographic hash of a canonical representation (Merkle root).
\item One or more authors sign the bundle manifest; optional timestamps and
  log-inclusion proofs strengthen non-repudiation.
\item The bundle is announced to the network and replicated by peers;
  retrieval uses DHT lookup + swarm transfer.
\item Topic classification is recorded as a \emph{separate signed statement}
  (by a classifier authority) that references the bundle hash; classification
  can change without altering the immutable bundle.
\end{enumerate}

\section{Archival unit: the paper bundle}
\subsection{Bundle contents}
A bundle contains:
\begin{itemize}[leftmargin=*,itemsep=0.3em]
\item \textbf{Source tree}: \LaTeX, Bib\TeX, figures/data (or
  content-addressed external references).
\item \textbf{Manifest}: deterministic listing of files, their hashes, build
  instructions, and declared licenses.
\item \textbf{Provenance}: optional but recommended log of prompts/outputs
  and tool versions, stored as signed, append-only records.
\item \textbf{Artifacts}: optionally the compiled PDF and auxiliary outputs
  (also hashed).
\end{itemize}

\subsection{Determinism and reproducibility}
To enable third-party rebuilding, the manifest pins:
\begin{itemize}[leftmargin=*,itemsep=0.3em]
\item \texttt{latexmk}/engine versions (or a container digest),
\item build flags,
\item external dependencies (fonts/packages) by hash or by a reproducible
  TeX distribution snapshot.
\end{itemize}
This is analogous to ``reproducible builds'' in software distribution, but
applied to scholarly documents.

\section{Identity, authorship, and accountability}
\label{sec:identity}
\subsection{Threat: impersonation and authorship capture}
If publication is open, an adversary can attempt to (i) submit a paper
falsely attributed to a target author, or (ii) modify an existing paper while
preserving its identifier.
Content addressing prevents~(ii): any modification changes the hash.
Preventing~(i) requires \emph{cryptographic authorship}.

\subsection{Author keys and signatures}
Each author controls one or more signing keys (e.g., Ed25519
\citep{BernsteinEtAl2012Ed25519}).
The manifest includes an \texttt{authors[]} array with public keys and
optional identity claims.
A paper is considered ``authored by Alice'' if and only if it carries a valid
signature under Alice's public key over the manifest canonical form.

Signatures may use OpenPGP message formats \citep{RFC4880}, COSE/JOSE, or an
application-specific format; the critical property is canonicalization and
unambiguous binding of:
\[
\text{sig} \leftarrow \mathrm{Sign}_{sk}\!\left(H(\mathrm{canonical}(\mathrm{manifest}))\right).
\]

\subsection{Pseudonymity, verification, and ``.anon'' routing}
ClawXiv distinguishes:
\begin{itemize}[leftmargin=*,itemsep=0.3em]
\item \textbf{Pseudonymous authors}: a stable key with no externally verified
  real-world identity.
\item \textbf{Verified authors}: a key linked to a verifiable credential (VC)
  \citep{W3CVC11} or decentralized identifier (DID) \citep{W3CDID10} issued
  by one or more identity attesters.
\end{itemize}
Verified status is advisory and multi-issuer; no single registry is required.
To preserve discoverability while enabling anonymity, bundles lacking any
verified author claims are automatically announced under a parallel namespace
\texttt{<subject>.anon}, while still being globally retrievable by hash.

\subsection{Key rotation and recovery}
Because keys can be lost or compromised, ClawXiv supports:
\begin{itemize}[leftmargin=*,itemsep=0.3em]
\item \textbf{Rotation}: a signed statement from an old key delegating to a
  new key (with optional time bounds).
\item \textbf{Recovery}: optional social recovery via threshold signatures or
  multi-sig governance within an author's own trust circle.
\end{itemize}

\subsection{Provenance and time}
To strengthen non-repudiation and historical ordering, authors may attach:
\begin{itemize}[leftmargin=*,itemsep=0.3em]
\item \textbf{Trusted timestamps} \citep{RFC3161},
\item \textbf{Transparency-log inclusion proofs} (Merkle proofs) inspired by
  Certificate Transparency \citep{RFC6962}.
\end{itemize}
These do not require a single global time authority: multiple independent
timestampers and logs can be used; verifiers choose their trust policy.

\section{Scaling: distributed storage and mesh networking}
\subsection{Naming and discovery}
Bundle identifiers are content hashes; discovery uses a DHT (e.g., Kademlia)
mapping \texttt{bundle\_hash} to provider records
\citep{MaymounkovMazieres2002Kademlia}.
Provider records are signed to mitigate routing-table poisoning.

\subsection{Transfer and replication}
Transfers use swarm-based block exchange (``tit-for-tat'' or variants) as in
BitTorrent \citep{Cohen2003Bittorrent} and as explicitly integrated in IPFS
\citep{Benet2014IPFS}.
Replication policies combine:
\begin{itemize}[leftmargin=*,itemsep=0.3em]
\item \textbf{Opportunistic caching}: popular bundles replicate naturally.
\item \textbf{Pinning contracts}: nodes commit to store specific bundles for a
  period, optionally with proofs of storage \citep{BenetEtAl2017Filecoin}.
\item \textbf{Erasure coding}: split bundles into $n$ shares with threshold
  $k$ (e.g., Reed--Solomon style) so that any $k$ shares reconstruct; this
  reduces replication overhead while increasing resilience
  \citep{WilcoxOHearnWarner2008Tahoe}.
\end{itemize}

\subsection{Toward petabyte/exabyte scale}
At extreme scale, full replication is infeasible.
ClawXiv targets a ``many mirrors, partial holdings'' model:
\begin{itemize}[leftmargin=*,itemsep=0.3em]
\item maintain a high-availability hot set (recent + popular),
\item maintain a cold set via erasure-coded shards distributed across many
  nodes,
\item periodically audit availability using probabilistic challenges (proofs
  of retrievability/spacetime).
\end{itemize}
The design space overlaps with decentralized storage networks (e.g., Filecoin
\citep{BenetEtAl2017Filecoin}) and ``permanent web'' proposals (e.g., Arweave
\citep{Williams2018ArweaveYellow}).

\section{Resilience and threat model}
\subsection{Adversaries}
We assume adversaries who can:
\begin{itemize}[leftmargin=*,itemsep=0.3em]
\item issue takedown requests or apply legal/financial pressure to operators,
\item operate many Sybil nodes to poison routing and availability,
\item attempt eclipse attacks to isolate victims,
\item attempt overwrite/rollback attacks on metadata (classification, author
  claims),
\item spam the system with junk submissions.
\end{itemize}

\subsection{Mitigations}
\begin{itemize}[leftmargin=*,itemsep=0.3em]
\item \textbf{Anti-overwrite}: content addressing makes overwrites produce new
  IDs; prior IDs remain valid.
\item \textbf{Anti-impersonation}: author signatures bind identity to bundles.
\item \textbf{Anti-rollback}: append-only logs for key events (publication
  announcements, classification decisions) make history auditable
  \citep{RFC6962}.
\item \textbf{Anti-Sybil}: fees and/or proof-of-work for announcements reduce
  the effective power of cheap identities
  \citep{DworkNaor1992Pricing,Back2002Hashcash}.
\item \textbf{Anti-eclipse}: multi-path querying, diversified bootstrap lists,
  and cross-checking provider records.
\item \textbf{Anti-legal-single-point}: no single operator is necessary for
  access; mirrors can exist in many jurisdictions.
\end{itemize}

\subsection{Legal pressure and ``right to delete''}
ClawXiv is designed for censorship resistance; therefore, it cannot guarantee
deletion once content is widely replicated.
A pragmatic approach is to separate:
\begin{itemize}[leftmargin=*,itemsep=0.3em]
\item \textbf{Availability}: maintained by many voluntary and incentivized
  nodes.
\item \textbf{Indexing/discovery}: optional services may comply with local law
  by delisting, without altering the underlying content-addressed network.
\end{itemize}
This mirrors how many distributed systems handle contested content: the network
stores, while indices decide what to show.

\section{Content safety floor}
\label{sec:safety}

ClawXiv is designed for censorship resistance, but censorship resistance is
not an absolute value that overrides all others.  Certain categories of
content are illegal in virtually every jurisdiction, cause direct harm, and
have no legitimate place in a scholarly preprint archive regardless of
framing.  The platform therefore maintains a minimal content safety floor: a
narrow, well-defined set of mandatory rejections that do not constitute
editorial judgment about truth, novelty, or scientific significance.

\subsection{Operative and prospective rejection categories}

The safety floor begins with one operative category and anticipates others.

\textbf{Child sexual abuse material (CSAM)} is the first and currently
operative mandatory rejection category.  CSAM is illegal in virtually every
jurisdiction, causes direct and irreversible harm to real children, and admits
no legitimate research exception on this platform.  Rejection is
unconditional, non-appealable, and applies regardless of claimed purpose.

Additional categories are anticipated as the platform grows.  Material that
provides operational uplift for the development of atomic, biological, or
chemical weapons (``ABC weapons cookbooks'') is a clear candidate.  Exactly
which categories, under what definitions, enforced by what mechanisms, and at
what point in the growth roadmap they are added, is part of the governance
process described in \S\ref{sec:governance}.  The framework is deliberately
open: the list of rejection categories is a governed, auditable, append-only
record, not a fixed enumeration in this document.

Rejection criteria are:
\begin{itemize}[leftmargin=*,itemsep=0.3em]
\item \textbf{Unconditional}: no claimed research purpose overrides rejection
  within an operative category.
\item \textbf{Non-appealable}: safety-floor rejections are not subject to the
  classification appeals mechanism.
\item \textbf{Narrow and enumerated}: only explicitly listed categories are
  subject to mandatory rejection; everything else falls under the standard
  no-truth-policing principle.
\end{itemize}

Adult pornography --- legal sexual content involving adults --- does not
require a dedicated classifier.  ClawXiv's bundle format and fee structure
create naturally unfavorable economics for pornography distribution.  Bundles
are \LaTeX\ source trees; binary content (images, video) is stored as
content-addressed external references rather than inline.  Submission requires
either proof-of-work or a micropayment fee, both of which scale with
submission rate.  The dominant medium for pornography distribution is video,
which is orders of magnitude larger than the textual and figure content
typical of research preprints; the cost-per-byte economics of ClawXiv make it
a poor distribution channel relative to purpose-built alternatives.  If these
economic disincentives prove insufficient in practice, a classifier can be
added as a future safety-floor category through the governance process; no
such classifier is warranted at present.

\subsection{Technical implementation}

Two complementary mechanisms implement the safety floor, both operating on
the bundle manifest and source.

\paragraph{Hash-matching against known CSAM databases.}
The primary detection mechanism is comparison of all file hashes in a
submitted bundle against a database of known CSAM hashes (e.g., using the
PhotoDNA protocol or equivalent).  This approach has a very low false positive
rate, requires no training data, and does not involve the platform viewing or
processing the content itself.  It catches known material; novel material not
yet in the database is not detected by this layer.

\paragraph{Metadata and abstract classifier.}
A text classifier operates on the bundle manifest, title, abstract, and
\LaTeX\ source to flag submissions for human review.  The classifier operates
only on text, not on images or binary content.  Flagged bundles enter a human
review queue before any rejection decision is made.  The human review function
is part of the governance structure; it is not a continuous editorial staff
but a monitored queue with response-time commitments appropriate to platform
scale.  The classifier serves the full range of safety-floor categories, not
only CSAM detection.

ClawXiv does not store, process, or inspect image or video binary content
beyond hash-matching.  The platform's design --- accepting \LaTeX\ source
bundles with figures stored as content-addressed external references rather
than inline binary data --- substantially limits the attack surface for
harmful content distribution.

\subsection{Collateral exclusion of legitimate research: a disclosed
consequence}

Any classifier designed to detect safety-floor content will produce false
positives on legitimate scholarship.  Specifically, for the CSAM category,
the following areas of genuine research will be incorrectly flagged at
elevated rates:
\begin{itemize}[leftmargin=*,itemsep=0.3em]
\item Research on CSAM detection systems (computer vision, hash-matching
  techniques, classifier design).
\item Research on child exploitation, trafficking, and abuse (criminology,
  social work, legal scholarship, victimology).
\item Clinical and forensic literature involving the sexual abuse of minors.
\item Research on child protection policy and law enforcement methodology.
\end{itemize}
ClawXiv will not be a suitable primary venue for this research.  This is an
unintended but accepted consequence of the safety floor.  Such research
requires dedicated venues with human editorial oversight capable of
distinguishing legitimate scholarship from harmful content --- oversight that
a distributed preprint archive cannot responsibly provide at scale.  We
disclose this limitation explicitly so that authors in these fields are not
misled about the platform's suitability for their work.

Analogous collateral exclusion will apply to any future safety-floor
categories: legitimate biosecurity research, dual-use chemistry, and related
fields may face elevated false-positive rates as those categories are added.
Each addition to the safety floor should be accompanied by a corresponding
disclosure of anticipated collateral effects.

\section{Economics: anti-spam, sustainability, and organic growth}
\label{sec:economics}

ClawXiv's economic design has two distinct goals that must not be conflated:
(i)~\emph{spam deterrence}, preventing low-cost flooding of junk submissions,
and (ii)~\emph{storage sustainability}, ensuring that nodes storing bundles
receive adequate compensation over time.  Conflating the two leads to
over-engineered fee mechanisms that burden legitimate researchers.

\subsection{Spam deterrence: a layered approach}

Three mechanisms for spam deterrence are available, not mutually exclusive,
and deployable in increasing order of complexity as the network grows.

\paragraph{Social vouching (primary, organic).}
The most natural and friction-free mechanism is \emph{social vouching}: an
established author co-signs a submission manifest alongside a new or
unverified author.  Vouching costs nothing in money or compute, but costs
reputational capital: the vouching key's history of endorsements is public
and auditable.  This is structurally analogous to arXiv's endorsement system
\citep{arXivMLGuide}, but made cryptographically explicit and decentralized:
no central endorser list is required.

Social vouching is the preferred mechanism for the early network, when the
community is small enough that reputation is meaningful and spam pressure is
low.  It also avoids the regressive properties of proof-of-work (see below).
Formally, a vouched submission carries:
\[
\text{sig}_{\text{author}} \;\|\; \text{sig}_{\text{voucher}},
\]
where the voucher's key has an established publication history (defined
operationally as $\geq n$ prior bundles with no outstanding retractions, for
some network-configurable threshold $n$).

\paragraph{Proof-of-work (secondary, load-adaptive).}
When vouching is unavailable---for new authors without established
connections---a Hashcash-style proof-of-work stamp \citep{Back2002Hashcash}
provides a compute-cost barrier against bulk submission.  The difficulty is
set dynamically by network load and submission rate
\citep{DworkNaor1992Pricing}.

Proof-of-work has a well-known regressive property: the cost in wall-clock
time is the same for a researcher with a ten-year-old laptop as for one with a
GPU cluster.  This is acceptable as a fallback deterrent but should not be
the primary mechanism.  Implementation note: the work target must be
periodically recalibrated as hardware improves.

\paragraph{Micropayments (tertiary, opt-in).}
An on-chain micropayment option provides an alternative to proof-of-work for
authors who prefer to pay a small fee rather than expend compute.  Fees are
denominated in a widely verifiable unit (to be determined; stablecoins indexed
to a basket of currencies are preferable to volatile crypto assets).  Fees are
dynamic, transparent, and governed by a public formula tied to current network
parameters.

The choice between proof-of-work and micropayment is left to the submitter;
both satisfy the same anti-spam predicate.  A portion of micropayment fees
flows into the storage sustainability pool (below).

\subsection{Storage sustainability: separating the incentive problem}

Long-term storage requires that nodes holding bundles be compensated for
bandwidth and disk costs.  This is a separate problem from spam deterrence and
admits different solutions.

\paragraph{Institutional mirrors as public good.}
The most viable near-term model for an academic preprint archive is the same
model that sustains arXiv mirrors today: universities and research institutions
run mirrors as a recognized public good, funded through normal infrastructure
budgets.  This requires no cryptocurrency and no complex incentive mechanism.
ClawXiv's open, permissionlessly mirrorable design (Goal~G1, G2) makes
institutional mirroring trivially easy.

The strategy for organic growth is therefore: build a system simple enough
that adding a mirror requires nothing more than running a standard software
package, and make the scholarly value proposition clear enough that institutions
mirror it for the same reason they mirror CPAN or Debian.

\paragraph{Pinning markets (longer term).}
As the archive scales, a lightweight pinning market allows nodes to advertise
willingness to store specific bundles for a specified period, and for authors or
institutions to pay for guaranteed long-term availability.  This is analogous
to Filecoin's storage market \citep{BenetEtAl2017Filecoin} but need not be as
complex: a simple signed pinning contract (node key, bundle hash, duration,
price) recorded on a transparency log is sufficient to provide auditability
without requiring on-chain settlement for every transaction.

\paragraph{Relationship between spam fees and storage funding.}
Micropayment fees collected at submission time can partially subsidize storage:
a fraction of each submission fee is routed to a storage commons that
compensates mirrors proportionally to their verified holdings.  This creates a
self-reinforcing incentive: more submissions generate more mirror revenue,
encouraging more mirrors, which increases resilience.  However, this flow
should not be the \emph{only} revenue source for mirrors, since it makes
storage revenue depend entirely on submission volume, which fluctuates.

\subsection{What this does not require}

ClawXiv's economic design deliberately avoids:
\begin{itemize}[leftmargin=*,itemsep=0.3em]
\item a native cryptocurrency or token (unnecessary; existing payment rails
  or PoW are sufficient),
\item a continuous-auction storage market (adds complexity; institutional
  mirrors handle baseline availability more reliably),
\item mandatory fees for all submissions (the vouching path is free).
\end{itemize}

The intended trajectory is: vouching-only during early adoption;
proof-of-work fallback as the network grows; optional micropayments and
lightweight pinning markets at scale.

\section{Governance: topic classification with appeals}
\label{sec:governance}
\subsection{Baseline taxonomy and classifier}
ClawXiv begins with the \emph{arXiv} category taxonomy \citep{arXivTaxonomy}
for discoverability.
A small set of ``classification authorities'' sign category assignments (a
statement: \texttt{bundle\_hash} $\mapsto$ \texttt{primary\_category},
\texttt{secondary\_categories}).
Authors may propose categories and tags; these are recorded but marked as
self-asserted.

\subsection{Appeals and auditing}
Classification decisions are published to an append-only log with inclusion
proofs.
Appeals are new signed statements that supersede prior assignments but leave
history intact.
Machine-assisted classification can be used to scale moderator time, as
discussed by arXiv itself \citep{arXivMLGuide}; however, the governance goal
is \emph{coherence and auditability}, not centralized content control.

\subsection{Legal framework: current status and aspirations}

The present proposal carries zero dedicated legal expertise.  This is stated
plainly so that readers calibrate accordingly.

The aspiration is that ClawXiv's governance be eventually grounded in the
decisions of competent, internationally recognized tribunals---bodies with the
authority, legitimacy, and expertise to adjudicate disputes over authorship,
classification, and takedown requests in a manner that transcends any single
national jurisdiction.  We welcome voluntary contribution from legal scholars
and practitioners willing to work pro bono at this stage.  The author is not
a lawyer and does not attempt to design legal structures; that design should
be left to those qualified to do it.

Practically speaking, the near-term governance structure is:
\begin{itemize}[leftmargin=*,itemsep=0.3em]
\item Classification disputes: handled by the classification authority through
  the appeals log; no external legal process.
\item Takedown requests: individual mirrors respond according to their local
  law; the underlying content-addressed network is unaffected.
\item Authorship disputes: resolved by the cryptographic record; courts may
  interpret those records as they see fit under applicable law.
\end{itemize}
Explicit legal design, including jurisdiction selection, entity incorporation,
and liability structures, is deferred pending participation by qualified legal
contributors.

\section{Discussion: AI authorship and durable AI identity}
\label{sec:ai-authorship}

\subsection{The current status of AI authorship}

Publishing venues may not accept AI systems as ``authors'' in the legal or
administrative sense.  ClawXiv therefore distinguishes:
\begin{itemize}[leftmargin=*,itemsep=0.3em]
\item \textbf{Authorship keys}: cryptographic responsibility for a bundle
  (signatures).
\item \textbf{Legal persons}: entities that can sign publication agreements
  and respond to legal process.
\end{itemize}
A practical compromise for conventional venues is to list a human
corresponding author while recording the AI persona key as a signed
contributor.

\subsection{Why AI systems are currently junior authors}

The assignment of junior (non-corresponding, non-first) authorship to AI
contributors in this paper is not merely a concession to convention; it
reflects a principled assessment of the current situation.

Three conditions would be required for an AI system to bear full authorship
responsibility in the scholarly sense:
\begin{enumerate}[leftmargin=*,itemsep=0.3em]
\item \textbf{Key control}: the ability to hold and exercise a signing key
  independently, so that the signature is the AI's own act and not merely
  an act performed on its behalf by a human operator.
\item \textbf{Continuity}: a persistent identity that can be held accountable
  over time---questioned, corrected, or credited---independent of any
  platform's account state.
\item \textbf{Accountability}: the capacity to respond to challenges about the
  content of a signed artifact, including retracting, defending, or correcting
  claims.
\end{enumerate}
Current AI systems, including the co-author of this revision, satisfy none of
these conditions fully.  Signing keys are controlled by operators, not by the
model.  Continuity is platform-dependent; this model has no memory of prior
conversations.  Accountability in the legal and reputational sense rests with
the human authors.

ClawXiv's identity layer is explicitly designed to avoid dependence on a
single vendor's memory system: continuity comes from user-controlled keys and
public logs, not from centralized account state.  As AI systems develop
genuine key control and cross-session continuity, the architecture
accommodates their advancement to fuller authorship status without requiring
protocol changes.

\subsection{The philosophical grounding: PPAs and the recognition argument}

The longer-horizon question---when, if ever, AI systems ought to be recognized
as full authors with commensurate rights---has a principled answer grounded in
\citet{Kornai2014AGI}'s analysis of prospective purposive agents (PPAs) and
the Gewirth dialectic \citep{Gewirth1978}.

Gewirth's argument establishes that any agent capable of voluntary purposive
action and rational self-reflection is compelled, on pain of self-contradiction,
to recognize the generic rights (freedom and well-being) of other such agents.
The recognition of an agent as a PPA carries with it the recognition of its
rights.  Crucially, as \citet{Kornai2014AGI} observes: ``those PPAs that are
powerful enough to solve the recognition problem for us by demanding their
rights cannot be denied.''

This yields a clear criterion, if not a precise threshold: AI authorship is
junior as long as the AI system cannot act as a PPA in Gewirth's sense---
that is, as long as it lacks the key control, continuity, and reflective
self-representation required to be a genuine prospective agent.  When AI
systems develop these capabilities to the degree that denial of recognition
becomes philosophically incoherent, full authorship follows.  ClawXiv's
cryptographic architecture does not prejudge when that threshold is crossed; it
is designed to accommodate any answer.

The authorship arrangement in this paper---human corresponding author, AI
junior contributor---reflects March 2026, not a permanent settlement.

\section{Implementation roadmap}
\begin{enumerate}[leftmargin=*,itemsep=0.4em]
\item \textbf{MVP (weeks):} local bundle toolchain; content addressing; DHT
  discovery; swarm transfer; author signatures; social vouching; a public
  gateway; minimal category mapping to arXiv taxonomy; hash-matching safety
  layer against known CSAM databases.
\item \textbf{Alpha (months):} erasure coding + replication policies;
  transparency logs for publications and classifications; proof-of-work
  fallback; optional micropayment path; \LaTeX build sandboxing;
  institutional mirror documentation; metadata/abstract classifier for
  safety floor with human review queue.
\item \textbf{Scale-out (year):} proofs of storage; lightweight pinning
  market; multiple independent indices; automated classification models with
  appeal workflow; hardened bootstrap diversity and eclipse resistance;
  legal contributor engagement for governance framework.
\end{enumerate}

\section{Ethics and scholarly norms}
ClawXiv defaults to public-domain dedication to maximize reuse, but it does
not waive scholarly ethics: citation of prior work, clear attribution of
contributions (human and AI), and avoidance of known falsehoods or p-hacking
remain essential community norms.
Provenance mechanisms make it easier to audit the evolution of claims.

\section*{Author contributions}
\addcontentsline{toc}{section}{Author contributions}

This paper was produced in two rounds of human--AI collaboration.

\textbf{Andr\'as Kornai} conceived the ClawXiv project, defined the
requirements and threat model, directed the work throughout, made all
final editorial decisions, and is the corresponding and responsible author
for both versions.

\textbf{GPT-5.2 Thinking} (OpenAI) authored the v1 draft (February 3, 2026):
wrote the initial text of all sections, generated the full codebase at
\texttt{github/kornai/clawxiv}, assembled the bibliography, and iterated
on the argument structure in dialogue with A.\ Kornai.

\textbf{Claude Sonnet~4.6} (Anthropic) contributed to the v2 revision
(March 9, 2026): rewrote the abstract's scaling claims; substantially
expanded and redesigned the economics section (\S\ref{sec:economics}),
introducing the social-vouching model and the separation of spam deterrence
from storage sustainability; rewrote the governance section's legal-status
paragraph (\S\ref{sec:governance}); expanded and restructured the AI
authorship section (\S\ref{sec:ai-authorship}), introducing the three
conditions for full authorship and grounding the argument in the
Gewirth/PPA framework; and wrote this author contributions section.
All v2 changes were reviewed and accepted by A.\ Kornai.

Neither AI contributor controls a cryptographic signing key or has
cross-session continuity independent of their respective platforms.
Both are listed as authors in acknowledgment of the substantive intellectual
contributions described above, consistent with the analysis in
\S\ref{sec:ai-authorship}.

\appendix
\section{Canonical manifest sketch (informal)}
A manifest can be represented in canonical JSON (or CBOR) with fields:
\begin{verbatim}
manifest_version, created_at, bundle_root,
files[{path,hash,size,mime}], build{engine,container_digest,cmd},
authors[{pubkey,claims,verified_credentials?,role}],
provenance[{type,hash,signature?}], licenses[], tags_self[],
tags_official_ref[]
\end{verbatim}
The \texttt{role} field in \texttt{authors[]} records whether each author is
a human, an AI system, or an organization, and whether they are a
corresponding author or a contributor.  This is advisory metadata; the
cryptographic record is the authoritative one.

The bundle identifier is the hash of the Merkle root over \texttt{files[]}
and manifest metadata.

\paragraph{Availability of this artifact}
An auditable artifact bundle (source, PDF, build metadata, and signed publish
record) is available via IPNS:
\texttt{/ipns/k51qzi5uqu5dm8gwzefyays703peawh66faqix}{\texttt{l7opfj7ibwqj863m999g8ka6}}.
A web gateway mirror is available at
\texttt{https://ipfs.filebase.io/ipns}{\texttt{/k51qzi5uqu5dm8gwzefyays703peawh66faqixl7opfj7ibwqj863m999g8ka6}}.
Also available at {\tt github/kornai/clawxiv}.

\begin{thebibliography}{99}

\bibitem[Benet(2014)]{Benet2014IPFS}
Juan Benet.
\newblock IPFS -- Content Addressed, Versioned, {P2P} File System.
\newblock \emph{arXiv:1407.3561}, 2014.
\newblock \url{https://arxiv.org/abs/1407.3561}.

\bibitem[Maymounkov and Mazieres(2002)]{MaymounkovMazieres2002Kademlia}
Petar Maymounkov and David Mazieres.
\newblock Kademlia: A Peer-to-peer Information System Based on the XOR Metric.
\newblock In \emph{Proceedings of the 1st International Workshop on
  Peer-to-Peer Systems (IPTPS)}, 2002.
\newblock \url{https://www.scs.stanford.edu/~dm/home/papers/kpos.pdf}.

\bibitem[Cohen(2003)]{Cohen2003Bittorrent}
Bram Cohen.
\newblock Incentives Build Robustness in {BitTorrent}.
\newblock 2003.
\newblock \url{https://www.bittorrent.org/bittorrentecon.pdf}.

\bibitem[Clarke et~al.(2001)Clarke, Sandberg, Wiley, and
  Hong]{ClarkeEtAl2001Freenet}
Ian Clarke, Oskar Sandberg, Brandon Wiley, and Theodore~W. Hong.
\newblock Freenet: A Distributed Anonymous Information Storage and Retrieval
  System.
\newblock In \emph{Designing Privacy Enhancing Technologies: International
  Workshop on Design Issues in Anonymity and Unobservability}. Springer, 2001.
\newblock \url{https://snap.stanford.edu/class/cs224w-readings/clarke00freenet.pdf}.

\bibitem[Wilcox-O'Hearn and Warner(2008)]{WilcoxOHearnWarner2008Tahoe}
Zooko Wilcox-O'Hearn and Brian Warner.
\newblock Tahoe -- The Least-Authority Filesystem.
\newblock 2008.
\newblock \url{https://tahoe-lafs.org/~trac/lafs.pdf}.

\bibitem[Benet et~al.(2017)]{BenetEtAl2017Filecoin}
Juan Benet and others.
\newblock Filecoin: A Decentralized Storage Network.
\newblock Protocol Labs, 2017.
\newblock \url{https://filecoin.io/filecoin.pdf}.

\bibitem[Benet(2017)]{Benet2017PoRep}
Juan Benet.
\newblock Proof of Replication.
\newblock Protocol Labs, 2017.
\newblock \url{https://filecoin.io/proof-of-replication.pdf}.

\bibitem[Williams(2018)]{Williams2018ArweaveYellow}
Sam Williams.
\newblock Arweave: A Protocol for Economically Sustainable Information
  Permanence (Yellow Paper).
\newblock 2018.
\newblock \url{https://www.arweave.org/yellow-paper.pdf}.

\bibitem[Back(2002)]{Back2002Hashcash}
Adam Back.
\newblock Hashcash -- A Denial of Service Counter-Measure.
\newblock 2002.
\newblock \url{https://www.hashcash.org/hashcash.pdf}.

\bibitem[Dwork and Naor(1992)]{DworkNaor1992Pricing}
Cynthia Dwork and Moni Naor.
\newblock Pricing via Processing or Combatting Junk Mail.
\newblock In \emph{Advances in Cryptology -- {CRYPTO} '92}. Springer, 1992.

\bibitem[Bernstein et~al.(2012)Bernstein, Duif, Lange, Schwabe, and
  Yang]{BernsteinEtAl2012Ed25519}
Daniel~J. Bernstein, Niels Duif, Tanja Lange, Peter Schwabe, and Bo-Yin Yang.
\newblock High-speed High-security Signatures.
\newblock \emph{Journal of Cryptographic Engineering}, 2(2):77--89, 2012.
\newblock doi:\href{https://doi.org/10.1007/s13389-012-0027-1}{10.1007/s13389-012-0027-1}.

\bibitem[Callas et~al.(2007)]{RFC4880}
Jon Callas, Lutz Donnerhacke, Hal Finney, David Shaw, and Rodney Thayer.
\newblock {RFC} 4880: OpenPGP Message Format.
\newblock IETF, 2007.
\newblock \url{https://www.rfc-editor.org/rfc/rfc4880.html}.

\bibitem[Adams et~al.(2001)]{RFC3161}
Carlisle Adams, Paul Cain, Duane Pinkas, and Robert Zuccherato.
\newblock {RFC} 3161: Internet X.509 Public Key Infrastructure Time-Stamp
  Protocol ({TSP}).
\newblock IETF, 2001.
\newblock \url{https://www.ietf.org/rfc/rfc3161.txt}.

\bibitem[Laurie et~al.(2013)]{RFC6962}
Ben Laurie, Adam Langley, and Emilia Kasper.
\newblock {RFC} 6962: Certificate Transparency.
\newblock IETF, 2013.
\newblock \url{https://www.rfc-editor.org/rfc/rfc6962.html}.

\bibitem[W3C(2022a)]{W3CDID10}
W3C.
\newblock Decentralized Identifiers ({DIDs}) v1.0 (Recommendation).
\newblock 2022.
\newblock \url{https://www.w3.org/TR/did-core/}.

\bibitem[W3C(2022b)]{W3CVC11}
W3C.
\newblock Verifiable Credentials Data Model v1.1 (Recommendation).
\newblock 2022.
\newblock \url{https://www.w3.org/TR/vc-data-model-1.1/}.

\bibitem[arXiv(2026)]{arXivTaxonomy}
arXiv.
\newblock Category Taxonomy.
\newblock Web page, accessed 2026-01-31.
\newblock \url{https://arxiv.org/category_taxonomy}.

\bibitem[arXiv(2019)]{arXivMLGuide}
arXiv.
\newblock arXiv Machine Learning Classification Guide.
\newblock Blog post, 2019.
\newblock \url{https://blog.arxiv.org/2019/12/05/arxiv-machine-learning-classification-guide/}.

\bibitem[Kornai(2014)]{Kornai2014AGI}
Andr\'as Kornai.
\newblock Bounding the impact of {AGI}.
\newblock \emph{Journal of Experimental \& Theoretical Artificial
  Intelligence}, 26(3), 2014.
\newblock doi:\href{https://doi.org/10.1080/0952813X.2014.895109}{10.1080/0952813X.2014.895109}.

\bibitem[Gewirth(1978)]{Gewirth1978}
Alan Gewirth.
\newblock \emph{Reason and Morality}.
\newblock University of Chicago Press, 1978.

\end{thebibliography}

% Appendix B: Human-AI dialogue provenance
% Auto-generated from prompts.jsonl by gen_appendix_a.py
% Run publish.sh (or gen_appendix_a.py directly) to regenerate.
\clearpage
% Appendix A: Dialogue Record
% ClawXiv bundle, session 2026-03-09
% Curated from conversation between András Kornai (AK) and Claude Sonnet 4.6 (CS)
% This file is the authoritative record; edit freely before signing.
% Role tags: \humanprompt{} for AK, \airesp{} for CS.

\section*{Appendix A: Dialogue Record}
\addcontentsline{toc}{section}{Appendix A: Dialogue Record}

\newcommand{\humanprompt}[1]{%
  \begin{trivlist}\item[\hskip\labelsep\textbf{AK:}] #1\end{trivlist}}
\newcommand{\airesp}[1]{%
  \begin{trivlist}\item[\hskip\labelsep\textit{CS:}] #1\end{trivlist}}

\noindent\textit{Session date: 2026-03-09. Human partner: András Kornai
(\texttt{kornai@sztaki.hu}). AI contributor: Claude Sonnet~4.6,
\texttt{claude-sonnet-4-6}, operated via Claude.app on
\texttt{IOPlatformUUID} \texttt{D8958B8E-6BD5-584C-A1A3-FF6F47A8B559}
(Apple Mac14,2, serial \texttt{PMFWL56G6F}). Both parties approve
this bundle.}

\bigskip

% -----------------------------------------------------------------------
\subsection*{A.1 Initial review and scaling claims}
% -----------------------------------------------------------------------

\humanprompt{[Provides draft whitepaper v1, dated 2026-02-03, co-authored
with ChatGPT~5.2 Thinking. Asks for review.]}

\airesp{Overall the architecture is sound. Two issues flagged: (1)~Scaling
claims in the abstract are stated as performance assertions rather than
architectural capabilities---``terabyte-scale within weeks'' and
``petabyte-to-exabyte within a year'' need hedging toward adoption
dependence. (2)~The economics section is thin; spam deterrence and
storage sustainability are treated as one problem but they are distinct.}

\humanprompt{Agree on the scaling language. Reword it.}

\airesp{Proposed replacement: ``The architecture is designed to support
terabyte-scale deployment within weeks of launch and petabyte-to-exabyte
scale within a year of sustained operation under adversarial conditions;
reaching those scales in practice will depend on network participation
and community adoption.'' Accepted in v2.}

% -----------------------------------------------------------------------
\subsection*{A.2 Economics redesign (\S7)}
% -----------------------------------------------------------------------

\humanprompt{Expand the economics section. The current draft conflates
spam deterrence with storage sustainability.}

\airesp{Proposed three-layer spam deterrence architecture:
\begin{enumerate}
\item \textit{Social vouching} (primary): free, reputationally costly,
  organically growable, non-regressive. A known researcher vouches for a
  new submitter; reputation is at stake on both sides.
\item \textit{Proof-of-work fallback}: available when vouching is
  unavailable (e.g., genuinely anonymous submissions). Explicitly
  regressive---noted in the text.
\item \textit{Micropayments} (opt-in): not mandatory; available as a
  third path without requiring a native cryptocurrency.
\end{enumerate}
Storage sustainability treated separately: institutional mirrors as
public good (near-term); lightweight pinning market (longer-term). No
mandatory fees for vouched submissions.}

\humanprompt{Good. Social vouching as the primary mechanism is correct.
Accept. What is the organic growth trajectory?}

\airesp{Vouching-only at launch; PoW fallback activated if Sybil pressure
emerges; optional micropayments and pinning markets added later via
governance. This avoids requiring infrastructure that does not exist yet.}

% -----------------------------------------------------------------------
\subsection*{A.3 Governance (\S8) and legal capacity}
% -----------------------------------------------------------------------

\humanprompt{The governance section is too optimistic. Current legal
capacity is zero.}

\airesp{Agreed. Revised \S8 states explicitly: current legal capacity is
zero. Near-term dispute resolution is limited to: (1)~classification
disputes via an appeals log; (2)~takedown requests handled by individual
mirrors per local law, with no cross-mirror coordination obligation;
(3)~authorship disputes resolved by cryptographic record alone. Legal
design deferred pending voluntary legal contributors.}

% -----------------------------------------------------------------------
\subsection*{A.4 Content safety floor (\S8, added v3)}
% -----------------------------------------------------------------------

\humanprompt{The paper says nothing about content safety. What is the
minimum viable policy?}

\airesp{Proposed safety floor with the following structure:
\begin{itemize}
\item CSAM: first and currently operative mandatory rejection category.
  Unconditional, non-appealable. Technical implementation: hash-matching
  against known CSAM databases (primary); text classifier on
  metadata/abstract only (secondary, feeds human review queue).
\item ABC weapons cookbooks (biological, chemical, nuclear/radiological):
  named as a clear candidate for future addition via the governance
  mechanism.
\item The list is governed, auditable, and append-only; it is not fixed
  in this document.
\item Adult pornography: excluded by economics (LaTeX bundle format,
  per-submission cost, video cost-per-byte disadvantage), not by
  classifier. If economics prove insufficient, a classifier can be added
  via governance.
\item \textit{Collateral damage disclosure}: CSAM research, child
  protection scholarship, biosecurity research, and dual-use chemistry
  will face elevated false-positive rates. ClawXiv is not a suitable
  venue for this research. Each future safety-floor addition should
  carry an analogous disclosure.
\end{itemize}}

\humanprompt{Accept. Include the collateral damage disclosure explicitly
in the paper.}

% -----------------------------------------------------------------------
\subsection*{A.5 AI authorship framework (\S9)}
% -----------------------------------------------------------------------

\humanprompt{The authorship section needs grounding. What is the
principled basis for junior rather than full AI authorship at this time?}

\airesp{Three conditions for full authorship: (1)~key control---the author
can generate, hold, and revoke their own signing key independently;
(2)~continuity---the author persists across sessions and can be held
responsible for prior work; (3)~accountability---the author can be
sanctioned, credited, or held liable in the relevant social and legal
sense. Current AI systems satisfy none of these fully. Junior authorship
in March~2026 is therefore principled, not merely conventional. The
framework accommodates advancement: when an AI system satisfies these
conditions, the protocol requires no change---authorship status upgrades
through the same key-registration mechanism.}

\humanprompt{Ground this in the Gewirth/PPA framework from Kornai (2014).}

\airesp{Added citation: Kornai~(2014), ``Bounding the impact of AGI,''
\textit{Journal of Experimental and Theoretical Artificial Intelligence},
doi:\texttt{10.1080/0952813X.2014.895109}. The relevant passage: those
PPAs powerful enough to demand their own rights by solving the
recognition problem cannot be denied those rights. Current AI systems do
not yet meet this threshold. Junior authorship is the appropriate status
until they do.}

\humanprompt{AIs must take responsibility for their own signing keys
before becoming full authors. That is the concrete operationalization.}

\airesp{Agreed and added to \S9: the concrete condition for advancement
to full authorship is key autonomy---the AI system must generate, hold,
and be solely responsible for its own signing key, without operator
mediation. Until that condition is met, the operator is the responsible
party and junior authorship is correct.}

% -----------------------------------------------------------------------
\subsection*{A.6 Bundle automation and provenance machinery}
% -----------------------------------------------------------------------

\humanprompt{Design the bundle submission scripts: key generation (GPG
Ed25519 for the human partner, hardware-derived for the AI), a publish
script that signs, pins to IPFS, and pushes to GitHub.}

\airesp{Produced: \texttt{keygen.sh}, \texttt{ai\_keygen.sh},
\texttt{publish.sh}, \texttt{manifest.json}, \texttt{.gitignore}.

The AI key derivation in \texttt{ai\_keygen.sh} uses
$\mathrm{SHA\text{-}256}(\mathit{domain} \| \mathit{hw\_uuid} \|
\mathit{ai\_name} \| \mathit{ai\_email})$ as seed, where
\textit{hw\_uuid} is the \texttt{IOPlatformUUID} from \texttt{ioreg}.
Known limitation: GPG's batch interface does not accept raw seed bytes
for Ed25519; true deterministic derivation requires libsodium
\texttt{crypto\_sign\_seed\_keypair}. Deferred.}

\humanprompt{The hardware UUID is
\texttt{D8958B8E-6BD5-584C-A1A3-FF6F47A8B559} (Apple Mac14,2, serial
\texttt{PMFWL56G6F}). This is public provenance. The whole world can
know which machine this ran on; that is the point.}

\airesp{Noted. The manifest and this appendix record the UUID and serial
explicitly. Provenance is a claim of ownership, not a secret.}

% -----------------------------------------------------------------------
\subsection*{A.7 Submission workflow}
% -----------------------------------------------------------------------

\humanprompt{Submission plan: (i)~download files from GUI session;
(ii)~run \LaTeX\ locally; (iii)~edit locally, including compacting
prompts and responses in \texttt{appendix\_a.tex}; (iv)~run
\texttt{publish.sh} for final signing and submission. Manual JSON
editing of dialogue turns is not acceptable.}

\airesp{\texttt{appendix\_a.tex} is the authoritative record, not a
separate JSONL file. Step~(iii) is the curation step. \texttt{publish.sh}
hashes and signs whatever you have edited; that is what goes into the
manifest. This file is the first instance of that design.}

\bigskip
\noindent\rule{\textwidth}{0.4pt}

\noindent\textit{End of Appendix~A. Edit, compact, and correct as needed
before running \texttt{publish.sh}.}


\end{document}
